Recent advances in large language models (LLMs) have demonstrated promising capabilities for automated GPU kernel generation, yet existing benchmarks suffer from exploitable verification loopholes, synthetic operator combinations, and unrealistic evaluation conditions.
We introduce \textsc{KernelBench 2.0}, a comprehensive and \textbf{standardized} benchmark for evaluating LLM-based CUDA kernel generation that marks a pivotal shift from bottom-up operator collection to \textbf{top-down, model-driven construction}.
Tasks are dynamically instantiated from the configuration files of real production models to form a reproducible canonical dataset, ensuring every generated kernel is directly usable in modern inference engines.
Complementing this design, we significantly expand operator coverage to include linear attention, state-space models, mixture-of-experts routing, and quantization kernels, ultimately covering XX\% of architectures in HuggingFace Transformers and 100\% of models supported by vLLM and SGLang.
We introduce \textbf{enhanced validation} that combines property-based unit tests with end-to-end model inference checks, effectively preventing "hacking" exploits that plagued prior benchmarks.
Furthermore, our integrated \textbf{actionable profiling} uses real input workloads to provide hardware-level feedback (via NVIDIA Nsight Compute) that directly translates to production optimization opportunities.
\textsc{KernelBench 2.0} thus serves not just as a standardized leaderboard, but as a bridge between automated code generation and production deployment.